% Generated draft: requires mathematical and source review.
\documentclass[11pt]{article}
\usepackage[margin=25mm]{geometry}
\usepackage{fontspec}
\setmainfont{Latin Modern Roman}
\usepackage{amsmath,amssymb,mathtools,unicode-math}
\setmathfont{Latin Modern Math}
\usepackage{enumitem,needspace,xurl,hyperref}
\setlength{\parindent}{0pt}
\setlength{\parskip}{5pt}
\setlength{\emergencystretch}{3em}
\newcommand{\N}{\mathbb{N}}
\newcommand{\Z}{\mathbb{Z}}
\newcommand{\Q}{\mathbb{Q}}
\newcommand{\R}{\mathbb{R}}
\newcommand{\C}{\mathbb{C}}
\newcommand{\F}{\mathbb{F}}
\newcommand{\A}{\mathbb{A}}
\newcommand{\PP}{\mathbb{P}}
\newcommand{\E}{\mathbb{E}}
\newcommand{\eps}{\varepsilon}
\newcommand{\dd}{\,\mathrm{d}}
\begin{document}
\section*{UnsolvedMath ID: 10400126}
\textbf{Source problem code:} catalogue entry
\subsection{Definitions and mathematical statement}
In the intended moduli-of-flat-connections setting, construct the measure denoted $\tau_M(A)^{1/2}$ on the source component $\mathcal M_i$, from the square root of Reidemeister torsion, extending the nondegenerate case referenced but not included. If $d\mu_i(A)$ denotes this measure, the required identity, preserving the extract's indices, is
\[
e^{\pi\mathrm i\widetilde I_j/4}v_j
=\int_{\mathcal M_i}
 e^{\pi\mathrm i(-2I_A+h_A^0+h_A^1)/4}\,d\mu_i(A).
\]
The quantities $I_A,\widetilde I_j,v_j$, the relation between $i$ and $j$, and their normalization are not defined in the record. The notation $h_A^q$ usually indicates dimensions of twisted cohomology, but its exact coefficient system and convention are not supplied.

For nonacyclic flat data, Reidemeister torsion is naturally determinant-line data rather than automatically a positive scalar. Producing its square-root density/measure and specifying orientations, stabilizers, singular strata, and integrability is therefore part of the requested construction. No arbitrary square-root or phase convention is silently chosen.

The source reports Conjectures 7.7, 7.9, and 7.11 proved together for mapping tori of finite-order diffeomorphisms of surfaces of genus at least two, citing [6]. It separately reports Conjecture 7.7 proved for all Seifert fibered spaces in [168], using calculations [353,354] and analytic estimates. The latter is not asserted to prove Conjecture 7.11 for every Seifert manifold.
\subsection{Short English statement}
Construct a square-root Reidemeister-torsion measure on the intended flat-connection moduli component that satisfies the displayed phase-weighted integral identity. Preserve the source's special-case results and its undefined phase/index conventions.
\subsection{Sources}
\begin{itemize}
\item Ohtsuki, \emph{Problems on invariants of knots and 3-manifolds} (2002), Conjecture 7.11, PDF page 107, \url{https://msp.org/gtm/2002/04/gtm-2002-04-024s.pdf}; accessed 2026-07-29, PDF-font-structure extraction.
\item Source [6], simultaneous finite-order mapping-torus cases of Conjectures 7.7, 7.9, 7.11; [168], Seifert cases of Conjecture 7.7; [353,354], supporting calculations. Bibliography entries are not supplied. The research summary additionally mentions Witten and Freed without specific citations.
\item Local registry \texttt{data/unsolvedmath/problems.json}, permanent UnsolvedMath ID $10400126$, source code catalogue entry, \emph{Conjecture 7.11}. Dataset attribution: Ulam AI and the UnsolvedMath Contributors.
\end{itemize}
\end{document}
